\chapter{From Ordinary Symmetry to Grassmann Algebra}

\lectureoverview{Supersymmetry is a symmetry, so we begin by recalling what that word means in quantum mechanics. A physical transformation is represented on Hilbert space by a unitary operator; continuous transformations are generated infinitesimally, and their commutators encode how the transformations compose. A genuine symmetry preserves energy, which is equivalent to commuting with the Hamiltonian. Supersymmetry then changes one assumption shared by all familiar examples: its generators turn bosons into fermions and fermions into bosons. Constructing such generators leads naturally to creation and annihilation operators, anticommutators, and the finite Grassmann algebra in which odd quantities square to zero.}

\section{Symmetry Relates Different Configurations}

A symmetry is not merely a decorative regularity. It is an operation that relates one physical configuration to another. Left-right reflection, for example, relates the two hands. If the person and the surrounding laws possessed exact left-right symmetry, corresponding properties of the two hands would agree. More generally, objects or states connected by an exact symmetry must share the physical quantities that the symmetry preserves. In particle physics that frequently means equal masses; in atomic physics it means equal energies for states that lie in the same rotational multiplet.

Our aim is supersymmetry, but it is useful first to recall continuous symmetries such as rotations and translations. The question of whether supersymmetry is "continuous" in the ordinary sense is already awkward. Its infinitesimal parameters are not ordinary commuting real numbers. It is best regarded as a graded extension of continuous symmetry, and the Grassmann arithmetic introduced below makes that statement precise enough for our present purposes. \editorialnote{Modern terminology describes supersymmetry as a continuous graded symmetry. The qualification explains the lecture's warning that the usual continuous-versus-discrete question does not fit supersymmetry cleanly.}

\section{Unitary Operations on Hilbert Space}

A quantum system is represented by a state vector \(\ket{\psi}\). A physical symmetry operation is represented by an operator \(U\) acting on that state:
\begin{equation}
\ket{\psi}\longmapsto U\ket{\psi}.
\end{equation}
For a spatial rotation, the axis and angle determine \(U\). If a state contains particles at several positions, the transformed ket represents all of them at their rotated positions. A translation works similarly: it replaces the state by one in which the entire configuration has been displaced.

\begin{classroomqa}
\audiencequestion{The state vector belongs to an often infinite-dimensional Hilbert space. Are we rotating Hilbert space or ordinary three-dimensional space? \lecturetimestamp{00:04:20}}
\lecturerresponse{Both descriptions refer to the same operation at different levels. The unitary transformation acts on the Hilbert-space vector and represents the physical rotation of the configuration in ordinary space. It is not literally a three-dimensional rotation of the abstract Hilbert space; it is the representation there of a spatial rotation.}
\end{classroomqa}

The essential condition is unitarity:
\begin{equation}
U^\dagger U=UU^\dagger=\mathbf 1,
\qquad
U^{-1}=U^\dagger.
\end{equation}
This preserves inner products and therefore probabilities:
\begin{equation}
\braket{U\psi}{U\psi}
=\bra{\psi}U^\dagger U\ket{\psi}
=\braket{\psi}{\psi}.
\end{equation}

\begin{figure}[H]
\centering
\includegraphics[width=0.56\textwidth]{lecture_04_figure_02.png}
\lectureframe{00:05:56.420}
\caption{The early board transition from a transformed state to the unitarity condition. The ket and daggered product are visible, but the lecturer obscures their right-hand ends; the complete equations appear in the derivation below.}
\end{figure}

Strictly speaking, Wigner's theorem also permits antiunitary symmetries such as time reversal. The transformations connected continuously to the identity, which are the subject here, are represented unitarily. \editorialnote{The antiunitary possibility is a standard qualification; it does not alter the lecture's discussion of rotations, translations, or supersymmetry generators.}

\section{Infinitesimal Transformations and Generators}

Doing nothing is represented by \(U=\mathbf 1\), not \(U=0\). A sufficiently small transformation differs from the identity by a term linear in its small parameter. For a rotation through a small angle \(\epsilon\),
\begin{equation}
U(\epsilon)=\mathbf 1+i\epsilon L+O(\epsilon^2).
\end{equation}
Unitarity to first order requires
\begin{align}
U^\dagger U
&=\left(\mathbf 1-i\epsilon L^\dagger\right)
  \left(\mathbf 1+i\epsilon L\right)+O(\epsilon^2)\\
&=\mathbf 1+i\epsilon(L-L^\dagger)+O(\epsilon^2),
\end{align}
and hence
\begin{equation}
L^\dagger=L.
\end{equation}
The operator \(L\) is the infinitesimal generator.

Rotations require three generators, one for each independent axis:
\begin{equation}
L_x,\qquad L_y,\qquad L_z.
\end{equation}
A rotation about an arbitrary axis \(\widehat{\bm n}\) is generated by
\begin{equation}
L_{\widehat{\bm n}}=\widehat{\bm n}\cdot\bm L.
\end{equation}

\begin{classroomqa}
\audiencequestion{Can the generators for the \(x\), \(y\), and \(z\) axes be added to make a rotation about another axis? \lecturetimestamp{00:07:50}}
\lecturerresponse{Yes. Their linear combinations generate rotations about intermediate axes. Three independent components are enough to build every infinitesimal spatial rotation.}
\end{classroomqa}

For rotations, these generators are the components of angular momentum. The notation is useful more generally because every continuous symmetry is understood locally through its generators.

\section{The Commutator as a Multiplication Table}

The rotation generators obey
\begin{align}
[L_x,L_y]&=iL_z,&
[L_y,L_z]&=iL_x,&
[L_z,L_x]&=iL_y,
\end{align}
or, compactly,
\begin{equation}
[L_i,L_j]=i\epsilon_{ijk}L_k.
\end{equation}
The Levi-Civita symbol vanishes if two indices agree and is \(+1\) or \(-1\) for even or odd permutations of \(1,2,3\). The important structural fact is closure: the commutator of two generators is another generator.

Why should a commutator describe the geometry of rotations? Perform a small rotation about \(x\), then one about \(y\), undo the \(x\)-rotation, and finally undo the \(y\)-rotation. The operations do not cancel because rotations about different axes do not commute.

\begin{classroomqa}
\audiencequestion{Can this noncommutation be pictured through an airplane's roll, pitch, and yaw? \lecturetimestamp{00:12:36}}
\lecturerresponse{Yes. Successive rotations about body axes provide an immediate geometric demonstration that changing the order changes the final orientation. The algebra below is the infinitesimal version of that familiar fact.}
\end{classroomqa}

Let
\begin{equation}
U_x(\epsilon)=\mathbf 1+i\epsilon L_x,
\qquad
U_y(\delta)=\mathbf 1+i\delta L_y.
\end{equation}
Because operators act on kets from right to left, the stated sequence is represented by
\begin{equation}
\mathcal C
=U_y^{-1}(\delta)U_x^{-1}(\epsilon)
 U_y(\delta)U_x(\epsilon).
\end{equation}

\begin{classroomqa}
\audiencequestion{Why may the undoing operations be written as inverses rather than daggers? \lecturetimestamp{00:13:25}}
\lecturerresponse{Because the transformations are unitary: \(U^\dagger=U^{-1}\). For an infinitesimal transformation, taking the inverse simply reverses the sign of the first-order term, \(U_x^{-1}=\mathbf 1-i\epsilon L_x+O(\epsilon^2)\).}
\end{classroomqa}

Multiplying the four factors on the board became a lively exercise: all terms linear in \(\epsilon\) or \(\delta\) cancel, the room tracked the mixed terms, and a temporary factor-of-two confusion exposed how easy it is to lose operator order. The clean second-order result is
\begin{align}
\mathcal C
&=\mathbf 1+\epsilon\delta[L_x,L_y]
  +O(\epsilon^2\delta,\epsilon\delta^2)\\
&=\mathbf 1+i\epsilon\delta L_z
  +O(\epsilon^2\delta,\epsilon\delta^2).
\end{align}
Thus the closed sequence leaves a small rotation about \(z\) through an angle proportional to \(\epsilon\delta\). \editorialnote{The displayed multiplication fixes the order and coefficient after the exploratory classroom calculation. Reversing the convention for the sequence reverses the sign but not the conclusion that the residual transformation is governed by the commutator.}

\begin{classroomqa}
\audiencequestion{Would changing the order of the inverse operations make the product the identity? \lecturetimestamp{00:25:23}}
\lecturerresponse{If each transformation is placed immediately beside its own inverse, the pairs cancel. The nontrivial commutator sequence deliberately separates those pairs; its failure to collapse measures the order dependence.}
\end{classroomqa}

Longer products do become more complicated, but they contain no new principle.

\begin{classroomqa}
\audiencequestion{What happens when there are more than three generators or many small transformations? \lecturetimestamp{00:25:34}}
\lecturerresponse{There are more commutators and more bookkeeping, but knowing the closed algebra lets one determine how arbitrarily many infinitesimal transformations combine. Finite transformations can be assembled from many small ones.}
\end{classroomqa}

In this sense the commutator algebra is the multiplication table of a continuous group written near the identity.

\section{Translations Complete the Spatial Example}

Spatial translations are generated by momentum. With the sign convention used in the lecture,
\begin{equation}
\begin{aligned}
T_x(a)&=\mathbf 1+iaP_x+O(a^2),\\
T_y(b)&=\mathbf 1+ibP_y+O(b^2),
\end{aligned}
\end{equation}
where \(a\) and \(b\) are distances rather than angles.

Two translations commute. Displace a point in \(x\), then in \(y\), undo the \(x\)-translation, and undo the \(y\)-translation; the point returns to its origin:
\begin{equation}
[P_i,P_j]=0.
\end{equation}

\begin{classroomqa}
\audiencequestion{Does translating or rotating a physical object require forces to start and stop the motion? \lecturetimestamp{00:28:45}}
\lecturerresponse{That is not the operation being defined. A symmetry transformation compares the original state with the mathematical state in which the whole configuration has been displaced or rotated. It does not describe a dynamical trajectory, acceleration, or work done while moving the apparatus.}
\end{classroomqa}

A rotation and a translation do not commute. Translate a point, rotate it about the origin, translate back along the original axis, and undo the rotation; the point does not return to where it began. The mixed algebra is
\begin{equation}
[L_i,P_j]=i\epsilon_{ijk}P_k.
\end{equation}
This relation says that momentum itself transforms as a vector under rotations. The exact index exercise is less important here than the general lesson: transformations of Euclidean space are encoded in the commutators among their rotation and translation generators.

\section{A Symmetry Must Preserve Energy}

We now pass from the composition of transformations to dynamics. If \(U\) is genuinely a symmetry of a system, applying it cannot change the energy. Let \(\ket{E}\) be an energy eigenstate:
\begin{equation}
H\ket{E}=E\ket{E}.
\end{equation}
The symbol \(E\) on the right is a number; \(\ket{E}\) is a vector labeled by that eigenvalue.

\begin{figure}[H]
\centering
\begin{minipage}[t]{0.43\textwidth}
\centering
\includegraphics[width=\linewidth]{lecture_04_figure_03.png}
\lectureframe{00:34:38.239}
\end{minipage}
\hfill
\begin{minipage}[t]{0.50\textwidth}
\centering
\includegraphics[width=\linewidth]{lecture_04_figure_04.png}
\lectureframe{00:36:22.376}
\end{minipage}
\caption{The transformed ket \(U\ket{E}\) and the defining eigenvalue equation \(H\ket{E}=E\ket{E}\), isolated in successive board frames.}
\end{figure}

If \(U\) is a symmetry, then \(U\ket{E}\) has the same eigenvalue:
\begin{equation}
HU\ket{E}=EU\ket{E}.
\end{equation}
Since \(E\) is a scalar, it commutes with \(U\), and the original eigenvalue equation may be used:
\begin{align}
HU\ket{E}
&=EU\ket{E}\\
&=UE\ket{E}\\
&=UH\ket{E}.
\end{align}
Therefore
\begin{equation}
(HU-UH)\ket{E}=0.
\end{equation}

\begin{classroomqa}
\audiencequestion{Could the derivation be repeated from its first relation? \lecturetimestamp{00:35:55}}
\lecturerresponse{Start with \(H\ket{E}=E\ket{E}\). Symmetry says \(U\ket{E}\) has the same energy, so \(HU\ket{E}=EU\ket{E}\). Move the scalar \(E\) through \(U\), replace \(E\ket{E}\) by \(H\ket{E}\), and subtract. The result is \((HU-UH)\ket{E}=0\).}
\end{classroomqa}

\begin{figure}[H]
\centering
\includegraphics[width=0.70\textwidth]{lecture_04_figure_05.png}
\lectureframe{00:38:51.066}
\caption{The full board derivation. Curved guide marks connect the equal-energy statement to the substitution \(E\ket{E}=H\ket{E}\); the last line is \((HU-UH)\ket{E}=0\).}
\end{figure}

The argument holds for every energy eigenstate. Assuming those states form a complete basis, an operator that annihilates every one of them is the zero operator:
\begin{equation}
[H,U]=0.
\end{equation}
This is the criterion for a time-independent unitary transformation to be a symmetry of the Hamiltonian.

For an infinitesimal transformation,
\begin{equation}
U(\epsilon)=\mathbf 1+i\epsilon L_i+O(\epsilon^2),
\end{equation}
the identity already commutes with \(H\), so each generator must satisfy
\begin{equation}
[L_i,H]=0.
\end{equation}

\begin{figure}[H]
\centering
\includegraphics[width=0.47\textwidth]{lecture_04_figure_06.png}
\lectureframe{00:42:14.296}
\caption{The symmetry criterion condensed on the board: \([U,H]=0\), followed by its generator form, standardized from the partly soft lower label as \([L_i,H]=0\).}
\end{figure}

If \(L_i\) has no explicit time dependence, the Heisenberg equation gives
\begin{equation}
\frac{dL_i}{dt}
=i[H,L_i]=0
\qquad(\hbar=1).
\end{equation}
Thus symmetry generators are conserved charges. Spatial translation symmetry gives momentum conservation; rotation symmetry gives angular-momentum conservation. \editorialnote{The explicit-time-derivative qualification is standard and is implicit in the lecture's conservation-law statement.}

\section{Transform the Whole Physical System}

The phrase "rotation is a symmetry" can be misused if only part of the apparatus is transformed. A free spin can be rotated without changing its energy. Place it in a fixed magnetic field and rotate only the spin, however, and its Zeeman energy changes. That does not disprove rotational symmetry; it means the field was held fixed while the spin was rotated relative to it. Rotate the spin and field together, or use an external field that is itself invariant, and the energy is unchanged.

The same issue appears in a hydrogen atom. Translating only the electron relative to the proton changes their separation and therefore the Coulomb energy. Translation symmetry acts on the complete electron-proton system, moving both by the same amount.

\begin{classroomqa}
\audiencequestion{Must an external field itself possess the symmetry? \lecturetimestamp{00:43:09}}
\lecturerresponse{Either the fixed external field must already be invariant under the operation, or it must be transformed along with the rest of the system. A relative rotation between a spin and an unrotated field is a different physical configuration, not the symmetry transform of the whole setup.}
\end{classroomqa}

Rotational degeneracy in an atom gives a clean positive example. For a fixed angular momentum \(l\), rotations mix states with magnetic quantum numbers
\begin{equation}
m=-l,-l+1,\ldots,l.
\end{equation}
In the absence of an external magnetic or electric field, all members of the multiplet have the same energy:
\begin{equation}
E_{l,m}=E_l.
\end{equation}
The symmetry relates distinct states and forces their degeneracy.

\section{What Supersymmetry Changes}

Every familiar symmetry discussed so far preserves boson-versus-fermion character. A rotation takes an electron to another electron, not to a photon. Color rotations mix the three colors of a quark but do not turn the quark into a gluon. Ordinary internal and spacetime symmetries may mix components or quantum labels, yet their generators are bosonic and do not change fermion parity.

Supersymmetry is different:
\begin{equation}
Q\ket{B}\sim\ket{F},
\qquad
Q\ket{F}\sim\ket{B}.
\end{equation}
Its multiplets cross the boson-fermion divide. In an exact supersymmetric theory, partners in one multiplet have the same mass. This is precisely the pairing required by the previous lectures: matched bosonic and fermionic loops can cancel the large ultraviolet corrections to scalar masses. The construction is difficult to visualize because no ordinary spatial motion turns an integer-spin state into a half-integer-spin state. We therefore proceed algebraically.

\section{A Prototype Supercharge}

Let \(A,A^\dagger\) annihilate and create a boson, and let \(C,C^\dagger\) annihilate and create a fermion. The simplest Hermitian operator that interchanges their one-particle states is
\begin{equation}
Q=C^\dagger A+A^\dagger C.
\end{equation}
The first term removes a boson and creates a fermion; the second removes a fermion and creates a boson. If \(Q\) acts on a one-boson state, the \(A^\dagger C\) term vanishes because there is no fermion to annihilate, while \(C^\dagger A\) performs the conversion. On a one-fermion state the roles reverse.

\begin{classroomqa}
\audiencequestion{Does this construction create a single particle that is a superposition of a boson and a fermion? \lecturetimestamp{00:52:37}}
\lecturerresponse{\(Q\) is an operator, not a state. Acting on a pure bosonic state it produces a fermionic state, and acting on a pure fermionic state it produces a bosonic state. Acting linearly on an already general superposition produces another superposition, just as any linear operator does.}
\end{classroomqa}

The room next tests the vacuum. Since both annihilation operators kill it,
\begin{equation}
A\ket{0}=0,\qquad C\ket{0}=0,
\end{equation}
we have
\begin{equation}
Q\ket{0}=0.
\end{equation}
The result is the zero vector, not a second copy of the vacuum.

The prototype is intentionally schematic. A spin-zero boson can be mapped to either component of a spin-one-half fermion, so the generator itself needs a spinor index:
\begin{equation}
Q_\alpha\ket{B}\sim\ket{F_\alpha}.
\end{equation}
In the lecture's simplified two-component picture,
\begin{equation}
Q_\alpha\ket{F_\beta}\sim\delta_{\alpha\beta}\ket{B}.
\end{equation}

\begin{classroomqa}
\audiencequestion{How can the state-space dimensions match when a spin-zero boson has one spin component and a spin-half fermion has two? \lecturetimestamp{00:55:02}}
\lecturerresponse{A single unindexed \(Q\) is insufficient. There are spinor components \(Q_\alpha\), one capable of producing each fermion spin component. The supercharge itself carries half a unit of spin because it changes spin by half a unit.}
\end{classroomqa}

The \(Q_\alpha\) are fermionic, or more precisely Grassmann-odd, operators. Every term in them contains an odd number of fermion operators, which is why they reverse fermion parity. Ordinary measurable observables are even. One should distinguish an odd supercharge from its odd transformation parameter: a finite supersymmetry transformation is built from an even combination such as \(\epsilon^\alpha Q_\alpha+\bar\epsilon_{\dot\alpha}\bar Q^{\dot\alpha}\), where both \(\epsilon\) and \(Q\) are odd. \editorialnote{The lecture informally calls the generators Grassmann numbers. More precisely, the supercharges are odd operators, while the anticommuting parameters are Grassmann numbers; their product is even and can appear in a finite transformation.}

\section{Bosonic and Fermionic Operator Algebras}

For bosonic modes labeled by \(i,j\),
\begin{align}
[A_i,A_j]&=0,&
[A_i^\dagger,A_j^\dagger]&=0,&
[A_i,A_j^\dagger]&=\delta_{ij}.
\end{align}
The labels may denote momentum, position, spin, or any complete one-particle basis. Because creation operators commute,
\begin{equation}
A_i^\dagger A_j^\dagger\ket{0}
=A_j^\dagger A_i^\dagger\ket{0}.
\end{equation}
The two-boson state is symmetric under interchange.

Fermionic operators obey anticommutation relations:
\begin{align}
\{C_i,C_j\}&=0,&
\{C_i^\dagger,C_j^\dagger\}&=0,&
\{C_i,C_j^\dagger\}&=\delta_{ij},
\end{align}
where
\begin{equation}
\{X,Y\}=XY+YX.
\end{equation}
For a single mode,
\begin{equation}
\{C^\dagger,C^\dagger\}=2(C^\dagger)^2=0,
\end{equation}
so
\begin{equation}
(C^\dagger)^2=0.
\end{equation}
The operator is not zero; its square is. It can create one fermion in the mode, but a second application gives zero. This is Pauli exclusion in operator language.

For distinct labels,
\begin{equation}
C_i^\dagger C_j^\dagger\ket{0}
=-C_j^\dagger C_i^\dagger\ket{0},
\end{equation}
which is the antisymmetry of the two-fermion state. The algebraic replacement
\begin{equation}
\begin{gathered}
\text{bosons: commutators}\\
\Updownarrow\\
\text{fermions: anticommutators}
\end{gathered}
\end{equation}
encodes the statistical distinction.

Many bosons may occupy one mode. A large photon occupation can therefore behave as a classical electromagnetic field with an ordinary numerical amplitude. Fermions cannot accumulate in one mode, and an odd fermionic field does not acquire an ordinary commuting classical value. This motivates a different arithmetic for fermionic bookkeeping.

\section{Grassmann Numbers and Parity}

Ordinary numbers commute:
\begin{equation}
\alpha_i\alpha_j=\alpha_j\alpha_i.
\end{equation}
Grassmann generators anticommute:
\begin{equation}
\theta_i\theta_j=-\theta_j\theta_i.
\end{equation}
Setting \(i=j\), over a number system in which \(2\neq0\), gives
\begin{equation}
\theta_i^2=0.
\end{equation}
No ordinary numerical value such as six or seven can be assigned to \(\theta_i\). These symbols form an algebra designed to keep track of fermionic signs.

Ordinary coefficients commute with Grassmann generators:
\begin{equation}
\alpha\theta_i=\theta_i\alpha.
\end{equation}

\begin{classroomqa}
\audiencequestion{What kind of object results when an ordinary number multiplies a Grassmann number? \lecturetimestamp{01:12:32}}
\lecturerresponse{It remains Grassmann-odd. Multiplication by an ordinary coefficient changes its scale, not its parity. The span of the generators is a vector space over the ordinary numbers.}
\end{classroomqa}

The algebra has a \(\mathbb Z_2\) parity. Ordinary coefficients and products of an even number of generators are even; a generator or a product of an odd number is odd:
\begin{equation}
\begin{aligned}
\text{even}\cdot\text{even}&=\text{even},\\
\text{even}\cdot\text{odd}&=\text{odd},\\
\text{odd}\cdot\text{odd}&=\text{even}.
\end{aligned}
\end{equation}

\begin{classroomqa}
\audiencequestion{Is this analogous to the multiplication rules for real and imaginary numbers, or a generalization like the introduction of \(i\)? \lecturetimestamp{01:14:25}}
\lecturerresponse{It is a generalization of arithmetic, but the analogy should not be pushed far. Imaginary numbers still commute and satisfy \(i^2=-1\); Grassmann-odd quantities anticommute and square to zero. The shared even-odd multiplication pattern does not make the algebras the same.}
\end{classroomqa}

A product such as \(\theta_1\theta_2\) is even and commutes with every homogeneous element of the Grassmann algebra, but it is not an ordinary number. It is nonzero while
\begin{equation}
(\theta_1\theta_2)^2=0.
\end{equation}

\begin{classroomqa}
\audiencequestion{If \(\theta_1\theta_2\) is even, why is it not simply an ordinary number? \lecturetimestamp{01:17:36}}
\lecturerresponse{Even parity means that it commutes in the graded algebra; it does not mean that it is a real or complex scalar. Its nilpotence distinguishes it immediately from a nonzero ordinary number.}
\end{classroomqa}

\section{Functions of Grassmann Variables}

With one generator, every polynomial terminates after first order:
\begin{equation}
f(\theta)=\alpha+\beta\theta.
\end{equation}
All powers \(\theta^n\) with \(n\geq2\) vanish. With two generators the most general element is
\begin{equation}
f(\theta_1,\theta_2)
=\alpha+\beta_1\theta_1+\beta_2\theta_2
+\gamma\theta_1\theta_2.
\end{equation}
Writing the last monomial in reversed order changes its sign:
\begin{equation}
\theta_2\theta_1=-\theta_1\theta_2.
\end{equation}

\begin{classroomqa}
\audiencequestion{Does multiplying \(\theta_1\) by a coefficient \(\beta\) create a third independent Grassmann variable? \lecturetimestamp{01:18:39}}
\lecturerresponse{No. It creates another vector in the span of \(\theta_1\), not a new basis generator. With two independent generators, every odd linear term is \(\beta_1\theta_1+\beta_2\theta_2\).}
\end{classroomqa}

For three generators, the ordered basis is
\begin{equation}
1,\quad
\theta_1,\theta_2,\theta_3,\quad
\theta_1\theta_2,\theta_1\theta_3,\theta_2\theta_3,\quad
\theta_1\theta_2\theta_3.
\end{equation}
In general, \(n\) generators produce \(2^n\) ordered basis monomials. Any analytic function of finitely many Grassmann generators truncates to such a polynomial because repeated generators vanish. \editorialnote{The \(2^n\)-dimensional basis states explicitly the pattern the lecture constructs for one, two, and three generators.}

\section{Grassmann Differentiation}

For one variable define
\begin{equation}
\frac{\partial}{\partial\theta}\alpha=0,
\qquad
\frac{\partial}{\partial\theta}\theta=1,
\end{equation}
so that
\begin{equation}
\frac{\partial}{\partial\theta}
(\alpha+\beta\theta)=\beta.
\end{equation}
For two variables, choose left derivatives and keep the canonical order \(\theta_1\theta_2\):
\begin{align}
\frac{\partial^L f}{\partial\theta_1}
&=\beta_1+\gamma\theta_2,\\
\frac{\partial^L f}{\partial\theta_2}
&=\beta_2-\gamma\theta_1.
\end{align}
The minus sign in the second line arises because the odd derivative must pass through \(\theta_1\) before differentiating \(\theta_2\).

This is expressed by the graded product rule. If \(A\) has parity \(|A|\in\{0,1\}\), then
\begin{equation}
\frac{\partial^L}{\partial\theta}(AB)
=\left(\frac{\partial^L A}{\partial\theta}\right)B
+(-1)^{|A|}A
\left(\frac{\partial^L B}{\partial\theta}\right).
\end{equation}
The rule is ordinary in spirit but contains the sign required by anticommutation. Grassmann integration can also be defined and, because the function space is finite, reduces to a short table. The lecture postpones that table.

\section{The Exponential Puzzle, Resolved}

With a single Grassmann generator and ordinary coefficient \(a\),
\begin{equation}
e^{a\theta}
=1+a\theta,
\end{equation}
because every higher power vanishes. Two exponentials involving the same generator obey
\begin{align}
e^{a\theta}e^{b\theta}
&=(1+a\theta)(1+b\theta)\\
&=1+(a+b)\theta\\
&=e^{(a+b)\theta}.
\end{align}

The room then tries distinct generators. Direct multiplication gives
\begin{equation}
e^{\theta_1}e^{\theta_2}
=(1+\theta_1)(1+\theta_2)
=1+\theta_1+\theta_2+\theta_1\theta_2.
\end{equation}
At first the board expansion suggests that the square of \(\theta_1+\theta_2\) might supply the last term. An audience member catches the sign:
\begin{equation}
(\theta_1+\theta_2)^2
=\theta_1^2+\theta_1\theta_2
 +\theta_2\theta_1+\theta_2^2
=0.
\end{equation}

\begin{classroomqa}
\audiencequestion{Does the square term cancel because it contains \(\theta_1\theta_2+\theta_2\theta_1\)? \lecturetimestamp{01:24:00}}
\lecturerresponse{Yes. That correction changes the conclusion: \(e^{\theta_1+\theta_2}=1+\theta_1+\theta_2\), while the ordered product \(e^{\theta_1}e^{\theta_2}\) also contains \(\theta_1\theta_2\). The naive exponential addition law does not apply to distinct anticommuting exponents.}
\end{classroomqa}

There is no contradiction. Even ordinary matrices satisfy \(e^Ae^B=e^{A+B}\) only when \(A\) and \(B\) commute. Here
\begin{equation}
[\theta_1,\theta_2]
=\theta_1\theta_2-\theta_2\theta_1
=2\theta_1\theta_2.
\end{equation}
This commutator is even and commutes with both generators, so the Baker-Campbell-Hausdorff series terminates:
\begin{equation}
e^{\theta_1}e^{\theta_2}
=\exp\left(
\theta_1+\theta_2+\frac12[\theta_1,\theta_2]
\right)
=e^{\theta_1+\theta_2+\theta_1\theta_2}.
\end{equation}
Since the exponent on the right is nilpotent, it expands to exactly
\(1+\theta_1+\theta_2+\theta_1\theta_2\). \editorialnote{This BCH calculation supplies the resolution that the live lecture leaves as a question. The lecture correctly discovers the failure of the naive law but ends before identifying the missing commutator term.}

\begin{classroomqa}
\audiencequestion{What kind of number is \(1+\theta_1\), and are ordinary and Grassmann parts allowed to be added? \lecturetimestamp{01:26:16}}
\lecturerresponse{It is an inhomogeneous element of the Grassmann algebra, with an even part \(1\) and an odd part \(\theta_1\). Such sums are allowed, but they are neither ordinary scalars nor homogeneous odd elements.}
\end{classroomqa}

Near the end, another apparent paradox is tested:
\begin{equation}
(\theta_1+\theta_2)^2=0.
\end{equation}
This does not imply \(\theta_1\theta_2=0\), because the two cross terms cancel rather than vanish separately.

\begin{classroomqa}
\audiencequestion{If every odd Grassmann element squares to zero, does that mean the product of any two Grassmann generators is zero? \lecturetimestamp{01:32:19}}
\lecturerresponse{No. The equation \((\theta_1+\theta_2)^2=0\) says \(\theta_1\theta_2+\theta_2\theta_1=0\). It proves anticommutation, not that either ordered product vanishes. Distinct generators can have a nonzero product.}
\end{classroomqa}

\begin{classroomqa}
\audiencequestion{Are we assuming the ordinary power-series definition when we write these exponentials? \lecturetimestamp{01:34:42}}
\lecturerresponse{Yes, and that definition is consistent because every series terminates. What fails is not the exponential itself but the unjustified identity \(e^Ae^B=e^{A+B}\) for noncommuting \(A\) and \(B\). The commutator correction restores the proper composition law.}
\end{classroomqa}

This exploratory ending is appropriate. Ordinary symmetries led us from unitary transformations to generators and commutators; supersymmetry led us from odd generators to anticommuting variables. In both cases, the multiplication law is the central object. For rotations it appears in \([L_i,L_j]\); for Grassmann exponentials it appears in the commutator term that prevents naive addition. The algebra is unfamiliar, but it is coherent, finite, and exactly suited to the fermionic transformations that supersymmetry requires.
